\clearpage\pagestyle{otMain}\setlength{\otBlocEspacement}{2.5mm}\fontsize{11}{15}\selectfont\corrigesChapitre{14}{Horaires et durées}
\begin{blocLivre}{cor-14-1}\begin{corrige}{1}{Lire deux horloges}\demarche{Lire la petite aiguille puis la grande.}
Première horloge : $\res{7\ \text h\ 15}$. Seconde : la grande aiguille sur le 7 indique 35 min, il est $\res{10\ \text h\ 35}$.
\label{fig-cor-14-1}\figureGeo[corrige]{hor-1}
\end{corrige}\end{blocLivre}
\begin{blocLivre}{cor-14-2}\begin{corrige}{2}{Placer les aiguilles}\demarche{Convertir l’heure de l’après-midi sur le cadran.}
Sur le cadran, 16 h se lit 4 h. La grande aiguille pointe le 9 (45 min) et la petite approche du 5 : le dessin figure en vert sur le corrigé.
\label{fig-cor-14-2}\figureGeo[corrige]{hor-2}
\end{corrige}\end{blocLivre}
\begin{blocLivre}{cor-14-3}\begin{corrige}{3}{Matin ou après-midi}\demarche{Ajouter ou retirer 12 heures.}
$15$ h $=\res{3\ \text h\ \text{de l’après-midi}}$ ; $20$ h $30=\res{8\ \text h\ 30\ \text{du soir}}$. « 5 h du soir » s’écrit $\res{17\ \text h}$.
\end{corrige}\end{blocLivre}
\begin{blocLivre}{cor-14-4}\begin{corrige}{4}{Convertir}\demarche{Passer par 60 minutes pour une heure.}
$2\times60=\res{120}$ min ; $60+30=\res{90}$ min ; $180=3\times60$, soit $\res3$ h.
\end{corrige}\end{blocLivre}
\begin{blocLivre}{cor-14-5}\begin{corrige}{5}{Dans la même heure}\demarche{Soustraire les minutes.}
Les heures sont les mêmes : $45-10=\res{35}$ min.
\end{corrige}\end{blocLivre}
\begin{blocLivre}{cor-14-6}\begin{corrige}{6}{Franchir l’heure}\demarche{Faire étape à l’heure ronde.}
De 9 h 40 à 10 h : 20 min. De 10 h à 10 h 15 : 15 min. Total : $20+15=\res{35}$ min.
\end{corrige}\end{blocLivre}
\begin{blocLivre}{cor-14-7}\begin{corrige}{7}{L’heure de fin}\demarche{Compléter d’abord jusqu’à l’heure ronde.}
$10$ h $20+40$ min $=11$ h, puis 5 min : fin à $\res{11\ \text h\ 05}$.
\end{corrige}\end{blocLivre}
\begin{blocLivre}{cor-14-8}\begin{corrige}{8}{L’heure de départ}\demarche{Reculer par étapes depuis la fin.}
On recule : 5 min pour revenir à 12 h, puis 25 min : départ à $\res{11\ \text h\ 35}$. Contrôle : $11$ h $35+30$ min $=12$ h 05.
\end{corrige}\end{blocLivre}
\begin{blocLivre}{cor-14-9}\begin{corrige}{9}{Des minutes aux heures}\demarche{Extraire les paquets de 60 minutes.}
$95=60+35$ : $\res{1\ \text h\ 35}$. Et $2\times60+5=\res{125}$ min.
\end{corrige}\end{blocLivre}
\begin{blocLivre}{cor-14-10}\begin{corrige}{10}{La plus longue durée}\demarche{Tout convertir en minutes.}
$1$ h $15=75$ min et $75<80$ : $\res{\text{Bob}}$ écoute 5 min de plus.
\end{corrige}\end{blocLivre}
\begin{blocLivre}{cor-14-11}\begin{corrige}{11}{La séance de cinéma}\demarche{Ajouter l’heure entière puis les minutes par étapes.}
$14$ h $50+1$ h $=15$ h 50, puis 10 min pour atteindre 16 h, puis 35 min : fin à $\res{16\ \text h\ 35}$.
\end{corrige}\end{blocLivre}
\begin{blocLivre}{cor-14-12}\begin{corrige}{12}{La journée de classe}\demarche{Calculer chaque demi-journée puis additionner.}
Matin : de 8 h 30 à 11 h 30, 3 h, puis 15 min : 3 h 15. Après-midi : $\res{3\ \text h}$ exactement. Total : $3$ h $15+3$ h $=\res{6\ \text h\ 15}$ de classe.
\label{fig-cor-14-12}\figureGeo[corrige]{hor-12}
\end{corrige}\end{blocLivre}
\begin{blocLivre}{cor-14-13}\begin{corrige}{13}{Le train de 9 h 48}\demarche{Découper le trajet aux heures rondes.}
De 9 h 48 à 10 h : 12 min. De 10 h à 11 h : 1 h. De 11 h à 11 h 23 : 23 min. Durée : $12+60+23=95$ min, soit $\res{1\ \text h\ 35}$.
\end{corrige}\end{blocLivre}
\begin{blocLivre}{cor-14-14}\begin{corrige}{14}{Les départs du bus}\demarche{Compter les intervalles puis ajouter le premier départ.}
Départs : 7 h, 7 h 20, 7 h 40, $\res{8\ \text h}$. De 7 h à 9 h, cela fait $120:20=6$ intervalles, donc $6+1=\res7$ départs : on compte le départ de 7 h.
\end{corrige}\end{blocLivre}
\begin{blocLivre}{cor-14-15}\begin{corrige}{15}{Le gâteau au four}\demarche{Comparer la durée réelle à la durée demandée.}
De 17 h 40 à 18 h 30 : $20+30=50$ min. Il fallait 55 min : il est sorti $\res{5\ \text{min trop tôt}}$. L’heure correcte de sortie était $\res{18\ \text h\ 35}$.
\end{corrige}\end{blocLivre}
